% ============================================================
%  Part VIII — Word order
% ============================================================
\gpart{VIII}{Word Order}

% ------------------------------------------------------------
\gsec{The verb-second rule}

In a German main clause the \textbf{conjugated verb is the second element}. Not
the second word — the second unit of meaning.

\begin{gtbl}{colspec={X[1.1,l] X[0.65,c] X[1.25,l]}}
  Position 1 & Verb & The rest \\
  Ich          & fahre & morgen nach Berlin. \\
  Morgen       & fahre & ich nach Berlin. \\
  Nach Berlin  & fahre & ich morgen. \\
  Weil es billig ist, & fahre & ich mit dem Bus. \\
\end{gtbl}

\begin{gnote}
Whatever you put in position 1 gets the emphasis, and the subject then slides in
straight after the verb. A whole subordinate clause counts as a single element,
which is why \textit{fahre} still comes second in the last example.
\end{gnote}

\tcap{The sentence bracket}

Everything verbal that is not the conjugated verb is parked at the end.

\begin{gtbl}{colspec={X[0.8,l] X[0.6,c] X[0.85,l] X[0.85,l]}}
  Pos.\,1 & Verb & Middle field & End \\
  Ich & habe & gestern ein Buch & gekauft. \\
  Ich & will & morgen früh & aufstehen. \\
  Er  & steht & jeden Tag um sechs & auf. \\
  Das Haus & wird & nächstes Jahr & gebaut. \\
\end{gtbl}

\tcap{Verb at the very end}

\begin{flattbl}{colspec={X[1,l] X[1.8,l]}}
  After a subordinating conjunction & \dots, weil ich müde \textbf{bin}. \\
  In a relative clause & Der Mann, der dort \textbf{steht}\dots \\
  In an indirect question & Ich weiß nicht, wo er \textbf{wohnt}. \\
\end{flattbl}

% ------------------------------------------------------------
\gsec{Order in the middle field}

\tcap{TeKaMoLo}

\begin{gtbl}{colspec={X[0.75,l] X[0.85,l] X[1.4,l]}}
  Slot & Asks & Example \\
  \textbf{Te}mporal & wann?  & heute, um acht, letzte Woche \\
  \textbf{Ka}usal   & warum? & wegen des Wetters, deshalb \\
  \textbf{Mo}dal    & wie?   & mit dem Bus, gern, schnell \\
  \textbf{Lo}kal    & wo? wohin? & in die Stadt, nach Hause \\
\end{gtbl}

\begin{flattbl}{colspec={X[1,l]}}
  Ich fahre \textbf{heute} \textbf{wegen des Wetters} \textbf{mit dem Bus} \textbf{in die Stadt}. \\
\end{flattbl}

\tcap{Objects}

\begin{gtbl}{colspec={X[1.2,l] X[1.8,l]}}
  Rule & Example \\
  Dativ before Akkusativ \eng{two nouns} & Ich gebe \textbf{dem Kind den Ball}. \\
  Pronoun before noun & Ich gebe \textbf{ihm} den Ball. \\
  Akkusativ before Dativ \eng{two pronouns} & Ich gebe \textbf{ihn ihm}. \\
\end{gtbl}

\begin{gnote}
The single rule underneath all three lines: \textbf{pronouns come early, and
known information comes before new information}. Two pronouns is the only case
where accusative jumps ahead of dative.
\end{gnote}

\tcap{Negation}

\begin{flattbl}{colspec={X[0.85,l] X[2.15,l]}}
  kein  & negates a noun with \textit{ein} or no article: \textit{Ich habe \textbf{kein} Auto.} \\
  nicht & negates everything else \\
  nicht & goes \textbf{before} the word it targets: \textit{nicht heute}, \textit{nicht schön} \\
  nicht & goes \textbf{late} when it negates the whole sentence, but still before the verb at the end \\
\end{flattbl}

% ------------------------------------------------------------
\gsec{Conjunctions}

Three families, three different effects on the verb. This is the single most
useful table in the chapter.

\tcap{Coordinating — position 0, word order untouched}

\begin{gtbl}{colspec={X[0.75,l] X[0.9,l] X[1.35,l]}}
  Word & Meaning & Example \\
  und     & and       & Ich komme \textbf{und} du bleibst. \\
  aber    & but       & Ich will, \textbf{aber} ich kann nicht. \\
  oder    & or        & Wir gehen \textbf{oder} wir bleiben. \\
  denn    & because   & Ich bleibe, \textbf{denn} ich bin müde. \\
  sondern & but rather & Nicht heute, \textbf{sondern} morgen. \\
\end{gtbl}

\tcap{Subordinating — verb to the end}

\begin{gtbl}{colspec={X[0.8,l] X[0.85,l] X[0.8,l] X[0.85,l]}}
  weil & because & obwohl & although \\
  dass & that & damit & so that \\
  wenn & if, whenever & bevor & before \\
  als & when \eng{one past event} & nachdem & after \\
  ob & whether & während & while \\
  falls & in case & seit(dem) & since \\
  bis & until & sobald & as soon as \\
  da & since \eng{reason} & solange & as long as \\
\end{gtbl}

\tcap{Adverbial — they occupy position 1, so the verb follows}

\begin{gtbl}{colspec={X[0.8,l] X[0.85,l] X[1.35,l]}}
  Word & Meaning & Example \\
  deshalb   & therefore & \textbf{Deshalb} \textbf{bleibe} ich zu Hause. \\
  deswegen  & for that reason & \textbf{Deswegen} \textbf{kam} er nicht. \\
  trotzdem  & nevertheless & \textbf{Trotzdem} \textbf{gehe} ich hin. \\
  dann      & then & \textbf{Dann} \textbf{essen} wir. \\
  außerdem  & besides & \textbf{Außerdem} \textbf{ist} es teuer. \\
  sonst     & otherwise & \textbf{Sonst} \textbf{wird} es zu spät. \\
\end{gtbl}

\begin{gnote}
\textit{denn} and \textit{weil} both mean \emph{because} but behave in opposite
ways: \textit{Ich bleibe, \textbf{denn} ich \textbf{bin} müde} versus
\textit{Ich bleibe, \textbf{weil} ich müde \textbf{bin}}.
\end{gnote}

\begin{gnote}
\textbf{als, wenn, wann.} \textit{als} for a single completed event in the past;
\textit{wenn} for \emph{if} and for repeated events; \textit{wann} only in
questions, direct or indirect.
\end{gnote}
